\chapter{Procesy rozgałęziające: losowy wzrost i wymieranie populacji}
\label{ch:branching}

\section{Od niezależnych osobników do zależnych pokoleń}
W procesie odnowy po każdej wymianie mechanizmu zaczyna się nowy
cykl. W procesie rozgałęziającym każdy osobnik produkuje losową
liczbę potomków, którzy w następnym pokoleniu rozmnażają się niezależnie.
Liczba osobników w kolejnym pokoleniu zależy od obecnej liczby,
ale cała genealogia powstaje z niezależnych lokalnych decyzji.
Zasadnicze pytanie brzmi: kiedy populacja wymiera, a kiedy ma
dodatnie prawdopodobieństwo przetrwania przez wszystkie pokolenia?

\begin{definicja}[Proces Galtona--Watsona]
Niech $\xi$ będzie zmienną losową o wartościach w $\N_0$,
z rozkładem $p_k=\Prob(\xi=k)$, $k\ge0$.
Niech $(\xi_{n,i})_{n\ge0,i\ge1}$ będą niezależnymi kopiami
$\xi$ oraz niech $Z_0=z\in\N_0$ będzie deterministyczne.
Definiujemy rekurencyjnie
\[
Z_{n+1}=\sum_{i=1}^{Z_n}\xi_{n,i},\qquad n\ge0,
\]
przyjmując sumę pustą za zero. Proces $(Z_n)_{n\ge0}$
nazywamy \emph{procesem rozgałęziającym Galtona--Watsona}.
Przyjmujemy filtrację
$\calF_n=\sigma(\xi_{k,i}:0\le k<n,\ i\ge1)$,
z trywialnym $\calF_0$.
\end{definicja}

Istnienie tablicy niezależnych zmiennych zapewnia
przeliczalny iloczyn przestrzeni z prawem $(p_k)_{k\ge0}$.
Indukcyjnie $Z_n$ jest skończoną zmienną całkowitoliczbową
i jest $\calF_n$-mierzalna: na zdarzeniu $\{Z_n=k\}$
suma definiująca $Z_{n+1}$ ma dokładnie $k$ składników,
a zdarzenia te tworzą przeliczalny podział przestrzeni.
Wiersz $(\xi_{n,i})_{i\ge1}$ jest niezależny od $\calF_n$.

\begin{uwaga}[Dlaczego występują dwa indeksy?]
Liczba $n$ identyfikuje pokolenie, natomiast $i$ rozróżnia
osobników w tym pokoleniu. Zmienna $\xi_{n,i}$ określa liczbę
dzieci $i$-tego osobnika w pokoleniu $n$. Niezależność jest
założeniem modelu: jeżeli wszystkie osobniki w danym pokoleniu
poddane są temu samemu losowemu wpływowi środowiska, zwykły
proces Galtona--Watsona może być nieadekwatny.
\end{uwaga}

\begin{center}
\begin{tikzpicture}[>=Latex,x=1cm,y=0.7cm]
\node[circle,draw,fill=blue!10,inner sep=2.5pt] (a) at (4,0) {};
\node[circle,draw,fill=blue!10,inner sep=2.5pt] (b) at (2,-1.8) {};
\node[circle,draw,fill=blue!10,inner sep=2.5pt] (c) at (4,-1.8) {};
\node[circle,draw,fill=blue!10,inner sep=2.5pt] (d) at (6,-1.8) {};
\node[circle,draw,fill=blue!10,inner sep=2.5pt] (e) at (1,-3.6) {};
\node[circle,draw,fill=blue!10,inner sep=2.5pt] (f) at (2,-3.6) {};
\node[circle,draw,fill=blue!10,inner sep=2.5pt] (g) at (5.3,-3.6) {};
\node[circle,draw,fill=blue!10,inner sep=2.5pt] (h) at (6.6,-3.6) {};
\foreach \x/\y in {a/b,a/c,a/d,b/e,b/f,d/g,d/h}{\draw[thick,blue!60!black] (\x)--(\y);}
\node[left] at (-.2,0) {$n=0$};
\node[left] at (-.2,-1.8) {$n=1$};
\node[left] at (-.2,-3.6) {$n=2$};
\node[right] at (7.3,0) {$Z_0=1$};
\node[right] at (7.3,-1.8) {$Z_1=3$};
\node[right] at (7.3,-3.6) {$Z_2=4$};
\end{tikzpicture}
\par\smallskip\small Przykładowe drzewo genealogiczne.
Środkowy osobnik pierwszego pokolenia nie ma potomstwa.
Proces zlicza liczby osobników, nie zapamiętuje pełnego drzewa.
\end{center}

\section{Własność Markowa i funkcje tworzące}
\begin{definicja}[Funkcja tworząca potomstwa]
Dla zmiennej $\xi$ o wartościach w $\N_0$ jej funkcją tworzącą
prawdopodobieństwa na $[0,1]$ jest
$f(s)=\E[s^\xi]=\sum_{k=0}^\infty p_ks^k$.
Przyjmujemy $s^0=1$, także dla $s=0$.
Średnią liczbę potomków oznaczamy
$m=\E\xi\in[0,\infty]$, a gdy istnieje drugi moment,
wariancję przez $\sigma^2=\Var(\xi)$.
Szereg definiujący $f$ jest jednostajnie zbieżny na
$[0,1]$ przez oszacowanie wyrazów przez $p_k$;
w szczególności $f$ jest ciągła i niemalejąca.
\end{definicja}

\begin{twierdzenie}[Rozkład przejścia i składanie funkcji tworzących]
Proces Galtona--Watsona z dowolnym deterministycznym
$Z_0=z\in\N_0$ jest jednorodnym łańcuchem Markowa
na $\N_0$. Jeżeli $Z_0=1$, to dla
$0\le s\le1$ jego funkcja tworząca $g_n(s)=\E[s^{Z_n}]$
spełnia
\[
g_0(s)=s,\qquad g_{n+1}(s)=g_n(f(s)),
\qquad g_n=f^{\circ n}.
\]
Przy starcie od $Z_0=z$ zachodzi $\E[s^{Z_n}]=(f^{\circ n}(s))^z$.
\end{twierdzenie}
\begin{proof}
Dla $k,j\in\N_0$ niech $P(k,j)$ będzie masą w $j$
$k$-krotnego splotu prawa $\xi$, przy czym splot zerowy
jest masą jednostkową w zerze.
Na zdarzeniu $\{Z_n=k\}$ niezależność $n$-tego wiersza
od $\calF_n$ daje
$\Prob(Z_{n+1}=j\mid\calF_n)=P(k,j)$.
Formalnie mnożymy przez funkcję charakterystyczną
dowolnego zdarzenia $A\in\calF_n$ i sumujemy po $k$.
Otrzymujemy $\Prob(Z_{n+1}=j\mid\calF_n)=P(Z_n,j)$,
czyli jednorodną własność Markowa, również względem
filtracji naturalnej (przez własność wieży).
Analogicznie dla $s\in[0,1]$ niezależność składników
daje $\E[s^{Z_{n+1}}\mid\calF_n]=f(s)^{Z_n}$.
Biorąc wartość oczekiwaną, otrzymujemy
$g_{n+1}(s)=\E[f(s)^{Z_n}]=g_n(f(s))$.
Indukcja i $g_0(s)=s$ dają $g_n=f^{\circ n}$.
Dla dowolnego $z$ ta sama rekurencja zaczyna się
od $g_0(s)=s^z$, więc daje $(f^{\circ n}(s))^z$.
\end{proof}

\begin{stwierdzenie}[Własność rozgałęziania]
\label{prop:gw-branching}
Prawo całego procesu z $z$ przodkami jest prawem sumy
$z$ niezależnych procesów Galtona--Watsona z jednym
przodkiem. Ponadto warunkowo względem $\calF_n$,
na $\{Z_n=k\}$ przyszła populacja ma prawo sumy
$k$ takich niezależnych procesów, z czasem liczonym od $n$.
\end{stwierdzenie}
\begin{proof}
Na przestrzeni iloczynowej budujemy niezależne kopie
procesu z jednym przodkiem. Jeśli ich obecne liczebności
wynoszą $k_1,\ldots,k_z$, rozkład sumy następnych
liczebności jest splotem $k_1+\cdots+k_z$ kopii
prawa $\xi$, niezależnie od podziału tej sumy.
Warunkowanie względem wspólnej historii i następnie
względem historii samej sumy daje macierz przejścia
$P$ z poprzedniego dowodu.
Stan początkowy sumy wynosi $z$, więc indukcja
po liczbie chwil identyfikuje jej rozkłady
skończeniewymiarowe z rozkładami pierwotnego procesu.
Dla drugiej tezy niezależność wszystkich przyszłych
wierszy tablicy od $\calF_n$ pozwala powtórzyć
ten rachunek przy ustalonym $Z_n=k$. Tożsamości dla
cylindrów rozszerzamy argumentem $\pi$--$\lambda$
na $\sigma$-ciało produktowe przyszłych trajektorii.
\end{proof}

\begin{twierdzenie}[Średnia liczba osobników]
Niech $m=\E\xi<\infty$ i $Z_0=z\in\N_0$.
Wtedy dla każdego $n\in\N_0$
\[
\E[Z_{n+1}\mid\calF_n]=mZ_n,\qquad
\E Z_n=z\,m^n,
\]
z konwencją $m^0=1$ także przy $m=0$.
Jeżeli $m>0$, proces $(Z_n/m^n)_{n\ge0}$ jest
nieujemnym martyngałem (dla $z\ge0$).
\end{twierdzenie}
\begin{proof}
Warunkowo względem $\calF_n$, po ustaleniu $Z_n=k$,
wartość oczekiwana sumy $k$ niezależnych kopii $\xi$
wynosi $km$. Ponieważ $\E Z_n<\infty$ wynika
indukcyjnie z $\E Z_0=z$, otrzymujemy
$\E Z_{n+1}=m\E Z_n$ i formułę $zm^n$.
Dla $m>0$, dzieląc równość warunkową przez $m^{n+1}$,
dostajemy $\E[Z_{n+1}/m^{n+1}\mid\calF_n]=Z_n/m^n$.
Całkowalność każdego wyrazu wykazaliśmy wcześniej.
\end{proof}

\section{Prawdopodobieństwo wymarcia jako punkt stały}
\begin{definicja}[Wymarcie]
Dla procesu Galtona--Watsona z $Z_0=1$ definiujemy
$\mathrm{Ext}:=\{\exists n\ge0:Z_n=0\}$ oraz
$q:=\Prob(\mathrm{Ext})$.
Ponieważ stan $0$ jest pochłaniający,
zdarzenie wymarcia jest sumą rosnącą zdarzeń $\{Z_n=0\}$.
\end{definicja}

\begin{twierdzenie}[Najmniejszy punkt stały]
Niech $f(s)=\E[s^\xi]$ będzie funkcją tworzącą potomstwa
procesu Galtona--Watsona z jednym przodkiem.
Wtedy prawdopodobieństwo wymarcia $q$ jest
\emph{najmniejszym} rozwiązaniem równania $f(s)=s$
w przedziale $[0,1]$. Można je obliczać rekurencją
$q_0=0$, $q_{n+1}=f(q_n)$, przy czym
$q_n\uparrow q$.
\end{twierdzenie}
\begin{proof}
Z poprzedniego twierdzenia
$q_n=\Prob(Z_n=0)=g_n(0)=f^{\circ n}(0)$,
co daje rekurencję. Zdarzenia $\{Z_n=0\}$ rosną,
więc $q_n$ rośnie i ma granicę $q\in[0,1]$.
Równość $q=\Prob(\mathrm{Ext})$ wynika z ciągłości miary
prawdopodobieństwa z dołu.
Funkcja tworząca $f$ jest ciągła na $[0,1]$:
ciągłość z lewej w $1$ otrzymujemy ze zbieżności monotonicznej
$s^\xi\uparrow1$, a w pozostałych punktach ze zbieżności
zmajoryzowanej. Dlatego przechodząc do granicy w
$q_{n+1}=f(q_n)$, dostajemy $q=f(q)$.
Jeżeli $r\in[0,1]$ spełnia $f(r)=r$, to $q_0=0\le r$.
Monotoniczność $f$ i indukcja dają $q_n\le r$
dla każdego $n$, a więc $q\le r$. To dowodzi minimalności.
\end{proof}

\begin{twierdzenie}[Kryterium wymierania]
Niech $\xi$ będzie zmienną o wartościach w $\N_0$, z
$m=\E\xi\in[0,\infty]$ oraz $p_0=\Prob(\xi=0)>0$.
Dla procesu Galtona--Watsona z jednym przodkiem:
\begin{enumerate}
\item jeśli $m\le1$, to $q=1$;
\item jeśli $m>1$, to $0<q<1$.
\end{enumerate}
Jeżeli $p_0=0$, wymarcie jest niemożliwe i $q=0$;
przypadek $\xi\equiv1$ ma $m=1$, lecz również $q=0$.
\end{twierdzenie}
\begin{proof}
Funkcja $f$ jest wypukła na $[0,1]$, bo każda
funkcja $s\mapsto s^k$ jest wypukła dla $k\ge0$.
Jej pochodna lewostronna w $1$ wynosi $m$,
również gdy $m=\infty$. Istotnie,
\[
\frac{1-f(s)}{1-s}
=\sum_{k\ge1}p_k\sum_{j=0}^{k-1}s^j
\uparrow\sum_{k\ge1}kp_k=m
\quad(s\uparrow1)
\]
na mocy twierdzenia o zbieżności monotonicznej.
Jeżeli $m\le1$, wypukłość daje dla $0\le s<1$
\[
f(s)\ge f(1)+m(s-1)=1-m(1-s)\ge s.
\]
Różnica $f(s)-s$ jest ściśle dodatnia dla $m<1$.
Jeżeli $m=1$ i $p_0>0$, warunek średniej wymusza
$\Prob(\xi\ge2)>0$, a więc $f$ jest ściśle wypukła
na $(0,1)$ i $f(s)>s$ również w tym przypadku.
Zatem jedynym punktem stałym jest $1$, stąd $q=1$.
Jeżeli $m>1$, to
$f(s)-s<0$ dla $s<1$ dostatecznie bliskich $1$:
istotnie iloraz $(f(s)-1)/(s-1)$ dąży do $m>1$.
Jednocześnie $f(0)-0=p_0>0$.
Ciągłość zapewnia istnienie punktu stałego $q\in(0,1)$.
Najmniejszy punkt stały, będący prawdopodobieństwem wymarcia,
leży w $(0,1)$: mniejszy od $1$ na mocy istnienia punktu stałego,
a większy od zera, gdyż $\Prob(Z_1=0)=p_0>0$.
Ponieważ $m>1$ wymusza $\Prob(\xi\ge2)>0$,
funkcja $f(s)-s$ jest ściśle wypukła na $[0,1]$.
Nie może mieć trzech różnych zer (wartość w środkowym
byłaby ściśle mniejsza od zera). Zatem $q$ jest
jedynym punktem stałym w $[0,1)$.
Gdy $p_0=0$, z prawdopodobieństwem $1$ wszystkie
zmienne z przeliczalnej tablicy są dodatnie,
więc $Z_n\ge1$ dla wszystkich $n$ i $q=0$.
\end{proof}

\begin{wniosek}[Wymarcie przy wielu przodkach i po zadanym pokoleniu]
Przy starcie z $z$ przodkami prawdopodobieństwo
wymarcia wynosi $q^z$ (także dla $z=0$, z $q^0=1$).
Dla dowolnego $n$ warunkowe prawdopodobieństwo
ostatecznego wymarcia względem $\calF_n$ wynosi $q^{Z_n}$.
\end{wniosek}
\begin{proof}
Dla startu z $z$ mamy
$\Prob(Z_n=0)=(f^{\circ n}(0))^z$.
Przejście do granicy daje pierwszą tezę.
Druga wynika z własności rozgałęziania
(stwierdzenie~\ref{prop:gw-branching}) albo z własności
Markowa dla zdarzeń $\{Z_{n+k}=0\}$ i warunkowej
zbieżności monotonicznej przy $k\to\infty$.
\end{proof}

\begin{zadanie}[Dwoje dzieci albo brak]
Każdy osobnik ma dwoje dzieci z prawdopodobieństwem $2/3$
albo nie ma dzieci z prawdopodobieństwem $1/3$.
Obliczyć prawdopodobieństwo wymarcia i średnią liczbę
osobników w pokoleniu $n$.

\emph{Rozwiązanie.} Funkcja tworząca ma postać
$f(s)=\frac13+\frac23s^2$, a średnia potomstwa
$m=4/3>1$. Równanie $f(q)=q$ daje
$2q^2-3q+1=0$, czyli pierwiastki $q=1/2$ i $q=1$.
Najmniejszy to $q=1/2$. Pomimo tego, że średnia
liczba osobników rośnie jak $(4/3)^n$, prawdopodobieństwo
całkowitego wymarcia wynosi aż $1/2$.
\end{zadanie}


\begin{zadanie}[Prawdopodobieństwo wymarcia dla trzech możliwości]
Niech liczba potomków ma rozkład
$\Prob(\xi=0)=1/4$, $\Prob(\xi=1)=1/4$,
$\Prob(\xi=2)=1/2$. Obliczyć średnią liczbę potomków
oraz prawdopodobieństwo ostatecznego wymarcia populacji
rozpoczętej od dwóch osobników.

\emph{Rozwiązanie.} Średnia wynosi
$m=0\cdot\tfrac14+1\cdot\tfrac14+2\cdot\tfrac12=5/4>1$.
Funkcja tworząca to
$f(s)=\tfrac14+\tfrac14s+\tfrac12s^2$. Równanie
$f(s)=s$ jest równoważne $2s^2-3s+1=0$,
a więc ma pierwiastki $1/2$ i $1$. Dla jednego przodka
$q=1/2$. Dwie niezależne linie genealogiczne wymierają
jednocześnie z prawdopodobieństwem $q^2=1/4$.
Przykład pokazuje, dlaczego na ogół nie wolno zastąpić
liczby przodków przez zmianę średniej $m$.
\end{zadanie}


\begin{zadanie}[Wymarcie w skończonej liczbie pokoleń]
Niech $\Prob(\xi=0)=\Prob(\xi=2)=1/2$ i $Z_0=1$.
Obliczyć prawdopodobieństwa wymarcia do drugiego i trzeciego
pokolenia; wyjaśnić różnicę między nimi a ostatecznym wymarciem.

\emph{Rozwiązanie.} Funkcja tworząca wynosi
$f(s)=(1+s^2)/2$. W pierwszym pokoleniu
$q_1=f(0)=1/2$. W drugim $q_2=f(q_1)=5/8$,
a w trzecim $q_3=f(q_2)=\tfrac12(1+25/64)=89/128$.
To konkretne kolejne przybliżenia rosną do $q=1$,
gdyż $m=1$ i $p_0>0$. Zdarzenie wymarcia do pokolenia
$n$ jest zdarzeniem $\{Z_n=0\}$, a ciąg tych zdarzeń
rośnie: raz wymarła populacja już się nie odradza.
\end{zadanie}

\section{Wariancja, normalizacja i granica martyngału}
\begin{twierdzenie}[Wariancja pokolenia]
Załóżmy, że $Z_0=1$, $0<m=\E\xi<\infty$ oraz
$\sigma^2=\Var(\xi)<\infty$. Wtedy
\[
\Var(Z_n)=
\begin{cases}
\displaystyle \sigma^2m^{n-1}\frac{m^n-1}{m-1},&m\ne1,\ n\ge1,\\[6pt]
 n\sigma^2,&m=1,\ n\ge0.
\end{cases}
\]
Jeżeli $m>1$, martyngał $W_n=Z_n/m^n$ jest ograniczony
w $L^2$, zatem zbiega w $L^2$ i prawie na pewno do
pewnej zmiennej $W_\infty\ge0$ o $\E W_\infty=1$ oraz
\[
\Var(W_\infty)=\frac{\sigma^2}{m(m-1)}.
\]
\end{twierdzenie}
\begin{proof}
Z niezależności potomstwa, na zdarzeniu $Z_n=k$,
$\E[Z_{n+1}^2\mid\calF_n]=k\sigma^2+k^2m^2$.
Indukcyjnie $\E Z_n^2<\infty$, zaczynając od $Z_0=1$.
W szczególności
$\Var(Z_{n+1}\mid\calF_n)=\sigma^2Z_n$.
Wzór na wariancję całkowitą daje dla $v_n=\Var(Z_n)$
\[
v_{n+1}=\E[\sigma^2Z_n]
+\Var(mZ_n)=\sigma^2m^n+m^2v_n,
\quad v_0=0.
\]
Indukcja daje
$v_n=\sigma^2\sum_{j=0}^{n-1}m^{2(n-1-j)+j}
=\sigma^2m^{n-1}\sum_{j=0}^{n-1}m^j$,
co po zsumowaniu szeregu geometrycznego daje wzór.
Dla $n=0$ wariancja wynosi zero. Dla $m>1$ otrzymujemy
$\Var(W_n)=m^{-2n}v_n
=\frac{\sigma^2}{m(m-1)}(1-m^{-n})$.
Ponieważ $\E W_n=1$, drugie momenty są jednostajnie
ograniczone. Z twierdzenia o zbieżności martyngałów
ograniczonych w $L^2$ dostajemy granicę w $L^2$ i p.n.;
przejście do granicy w pierwszych i drugich momentach daje
deklarowane wzory.
\end{proof}

\begin{uwaga}[Wartość oczekiwana granicy nie zawsze zachowuje się sama]
Dla $m>1$ nieujemny martyngał $W_n$ ma granicę prawie na pewno
nawet bez założenia skończonej wariancji, ale z tego
\emph{nie wynika} automatycznie $\E W_\infty=1$.
Skończoność wariancji jest tutaj wygodnym dostatecznym
warunkiem jednostajnej całkowalności. Bardziej ogólny
warunek zachowania średniej w procesie nadkrytycznym
wiąże się z twierdzeniem Kestena--Stiguma i momentem
$\E[\xi\log^+\xi]$; ten głębszy wynik można omówić
przy bardziej zaawansowanych procesach rozgałęziających.
Poniższe twierdzenie wyjaśnia, co można już
wywnioskować z niezerowej granicy martyngału.
\end{uwaga}

\begin{twierdzenie}[Granica martyngału a wymarcie]
\label{thm:gw-limit-extinction}
Niech $Z_0=1$ oraz $1<m<\infty$.
Martyngał $W_n=Z_n/m^n$ ma skończoną granicę
$W$ prawie na pewno. Jeżeli $\Prob(W>0)>0$, to
\[
\Prob(W=0)=q,\qquad
\{W=0\}=\mathrm{Ext}\quad\text{z dokładnością do zbioru zerowego}.
\]
W szczególności zachodzi to przy skończonej wariancji
potomstwa lub, ogólniej, przy jednostajnej całkowalności
$(W_n)_{n\ge0}$. Na zdarzeniu przetrwania w tych
przypadkach $Z_n/m^n\to W>0$ prawie na pewno.
\end{twierdzenie}
\begin{proof}
Twierdzenie o zbieżności nieujemnych martyngałów
daje granicę $W\ge0$ skończoną prawie na pewno,
z $\E W\le1$ na mocy lematu Fatou.
Na zdarzeniu wymarcia $Z_n=0$ od pewnego indeksu,
więc $W=0$.

Z własności rozgałęziania
(stwierdzenie~\ref{prop:gw-branching}) warunkowo
względem $Z_1=k$ proces od pierwszego pokolenia
ma prawo sumy $k$ niezależnych procesów z jednym
przodkiem. Ponieważ $k$ jest skończone, przechodząc
do granicy prawie pewnych granic na przestrzeni
produktowej, otrzymujemy równość rozkładów
\[
W\overset{d}=\frac1m\sum_{j=1}^{\xi} W^{(j)},
\]
gdzie $W^{(j)}$ są niezależnymi kopiami $W$,
niezależnymi od $\xi$. Można ją także uzyskać
z rozkładów warunkowych dla każdego skończonego $k$
i zsumować po $k$.
Nieujemność składników daje dla $r=\Prob(W=0)$
równanie $r=f(r)$. Jeśli $\Prob(W>0)>0$, to $r<1$.
Dla $p_0>0$ jedynym punktem stałym poniżej $1$
jest $q$. Dla $p_0=0$ i $m>1$ mamy dla $0<s<1$
$f(s)<s$, ponieważ $s^k\le s$ dla $k\ge1$
i $\Prob(\xi\ge2)>0$; jedynym punktem stałym
poniżej $1$ jest wtedy $q=0$.
Stąd $r=q$. Z zawierania $\mathrm{Ext}\subseteq\{W=0\}$
z dokładnością do zbioru zerowego i równości miar
wynika równość zdarzeń prawie na pewno.
Jeśli martyngał jest jednostajnie całkowalny,
zbiega w $L^1$, a $\E W=1$, więc $r<1$.
\end{proof}

\begin{wniosek}[Liczebność na zdarzeniu przetrwania]
Jeżeli $m>1$ (także $m=\infty$), to na zdarzeniu
przetrwania proces Galtona--Watsona z jednym
przodkiem spełnia $Z_n\to\infty$ prawie na pewno.
Nie trzeba zakładać skończonego drugiego momentu.
\end{wniosek}
\begin{proof}
Gdy $p_0>0$, mamy $0<q<1$. Warunkowe
prawdopodobieństwa wymarcia są równe $q^{Z_n}$.
Z twierdzenia o zbieżności warunkowych wartości
oczekiwanych względem rosnących $\sigma$-ciał
zbiegają prawie na pewno do
$\chi_{\mathrm{Ext}}$, ponieważ zdarzenie
wymarcia jest mierzalne względem
$\sigma(\bigcup_n\calF_n)$.
Na zdarzeniu przetrwania otrzymujemy więc
$q^{Z_n}\to0$, czyli $Z_n\to\infty$.
Jeżeli $p_0=0$, wszystkie osobniki mają co najmniej
jedno dziecko prawie na pewno, więc $Z_n$ jest
niemalejący. Przy $m>1$ mamy $p_1<1$.
Dla ustalonych $k\ge1$ i $n\ge0$ prawdopodobieństwo
pozostawania na poziomie $k$ przez następne $l$
pokoleń, pod warunkiem $Z_n=k$, wynosi $p_1^{kl}$.
Przejście $l\to\infty$ i suma po $k,n$ wykluczają
ostateczne pozostawanie na skończonym poziomie.
\end{proof}

\section{Całkowita liczba urodzonych osobników}
\begin{twierdzenie}[Średnia całkowita liczba osobników]
Niech $Z_0=1$ i $m=\E\xi<1$.
Zdefiniujmy całkowitą liczebność drzewa genealogicznego
$H=\sum_{n=0}^\infty Z_n$, gdzie suma może przyjmować $\infty$.
Wtedy $H<\infty$ prawie na pewno i
$\E H=1/(1-m)$.
\end{twierdzenie}
\begin{proof}
Zmienne $Z_n$ są nieujemne. Twierdzenie Tonellego
i wzór $\E Z_n=m^n$ dają
\[
\E H=\sum_{n=0}^\infty\E Z_n=\sum_{n=0}^\infty m^n
=\frac{1}{1-m}<\infty.
\]
Gdyby $\Prob(H=\infty)>0$, wartość oczekiwana $H$
byłaby nieskończona. Stąd $H<\infty$ prawie na pewno.
\end{proof}

\begin{stwierdzenie}[Funkcja tworząca całkowitej liczebności]
Niech $Z_0=1$ i $\Prob(H<\infty)=1$.
Dla $s\in[0,1]$ funkcja $h(s)=\E[s^H]$ spełnia
\[
h(s)=s f(h(s)),
\]
gdzie $f$ jest funkcją tworzącą liczby potomków.
\end{stwierdzenie}
\begin{proof}
Warunkowo względem $\xi=Z_1=k$, z własności
rozgałęziania (stwierdzenie~\ref{prop:gw-branching})
proces od pokolenia pierwszego ma prawo sumy
$k$ niezależnych drzew z jednym przodkiem.
Sumując liczebności po wszystkich pokoleniach,
otrzymujemy równość rozkładów
$H\overset{d}=1+\sum_{j=1}^\xi H_j$, gdzie po prawej
stronie możemy użyć niezależnych kopii $H$,
niezależnych od $\xi$, na przestrzeni iloczynowej.
Warunkowanie względem $\xi$ daje
$\E[s^H\mid\xi]=s\,h(s)^\xi$ prawie na pewno
(z warunkowej własności rozgałęziania), skąd po
uśrednieniu $h(s)=sf(h(s))$.
\end{proof}

\begin{zadanie}[Potomstwo Poissona]
Niech liczba potomków każdego osobnika ma rozkład
$\mathrm{Poisson}(\lambda)$, $0<\lambda<1$.
Obliczyć średnią całkowitą liczebność i równanie
dla jej funkcji tworzącej.

\emph{Rozwiązanie.} Mamy
$f(s)=\exp(\lambda(s-1))$ i $m=\lambda$.
Stąd $\E H=1/(1-\lambda)$ oraz
$h(s)=s\exp(\lambda(h(s)-1))$.
To równanie pozwala numerycznie obliczać
$\Prob(H=k)$, bez budowania całego losowego drzewa.
\end{zadanie}

\section{Model rozgałęziania w czasie ciągłym: proces Yule'a}
\begin{definicja}[Proces czystych urodzin Yule'a]
Niech $\lambda>0$. Procesem Yule'a rozpoczynającym się
od jednego osobnika nazywamy minimalny łańcuch Markowa
w czasie ciągłym na $\{1,2,3,\ldots\}$ z intensywnościami
$q_{n,n+1}=\lambda n$, $q_{nn}=-\lambda n$ oraz
$q_{n,j}=0$ dla pozostałych $j$.
Interpretacja: każdy z $n$ obecnych osobników rodzi
kolejnego niezależnie z intensywnością $\lambda$.
\end{definicja}

\begin{twierdzenie}[Brak eksplozji i rozkład geometryczny]
Proces Yule'a rozpoczynający się od $Z_0=1$ nie wybucha
w skończonym czasie. Dla wszystkich $t\ge0$ i $n\ge1$
\[
\Prob(Z_t=n)=e^{-\lambda t}(1-e^{-\lambda t})^{n-1},
\qquad \E Z_t=e^{\lambda t}.
\]
\end{twierdzenie}
\begin{proof}
Po przejściu ze stanu $n$ proces musi przejść do $n+1$.
Czas oczekiwania w stanie $n$ ma rozkład
$\mathrm{Exp}(\lambda n)$; kolejne czasy są niezależne.
Suma $n$ pierwszych czasów oczekiwania jest czasem przejścia
od $1$ do $n+1$. Jej transformata Laplace'a dla $c>0$ to
$\prod_{k=1}^n (\lambda k)/(\lambda k+c)$.
Gdyby czas eksplozji $\zeta=\sum_{k\ge1}H_k$ był
skończony z dodatnim prawdopodobieństwem, to
$\E e^{-c\zeta}>0$. Tymczasem
$\sum_{k\ge1}\log(1+c/(\lambda k))=\infty$,
ponieważ dla dużych $k$ wyrazy są nie mniejsze od
$c/(2\lambda k)$, a szereg harmoniczny jest rozbieżny.
Zbieżność zmajoryzowana, z konwencją $e^{-c\infty}=0$,
pozwala przejść do granicy w transformatach sum
częściowych. Zatem $\E e^{-c\zeta}=0$, co wymusza
$\Prob(\zeta=\infty)=1$. To jest bezpośredni dowód
przypadku kryterium eksplozji potrzebnego tutaj.

Dla przejść rozkładu $p_n(t)=\Prob(Z_t=n)$
uzasadnimy najpierw użycie równań w przód, chociaż przestrzeń
stanów jest nieskończona. Dla każdego ustalonego $K\ge2$
zastąpmy stany $K+1,K+2,\ldots$ jednym pochłaniającym
stanem $\partial$, nie zmieniając przejść z $1,\ldots,K$.
Do pierwszego osiągnięcia $K+1$ oba procesy są identyczne,
a ponieważ proces Yule'a wykonuje jedynie skoki $n\to n+1$,
dla $j\le K$ ich prawdopodobieństwa przebywania w stanie $j$
są równe. Proces obcięty ma skończoną przestrzeń stanów;
możemy więc zastosować wcześniej udowodnione równania
Kołmogorowa w przód (twierdzenie~\ref{thm:kolmogorov-eq}),
a otrzymane równania dla $p_j$, $j\le K$,
nie zależą od $K$. Dowodzi to dla wszystkich $n\ge1$ równań
$p_1'(t)=-\lambda p_1(t)$,
$p_n'(t)=-n\lambda p_n(t)+(n-1)\lambda p_{n-1}(t)$ dla $n\ge2$,
z warunkami $p_1(0)=1$ i $p_n(0)=0$ dla $n\ge2$.
Pierwsze równanie daje $p_1(t)=e^{-\lambda t}$.
Podstawiając kandydat
$p_n(t)=a(t)(1-a(t))^{n-1}$, gdzie $a(t)=e^{-\lambda t}$,
i korzystając z $a'=-\lambda a$, otrzymujemy po różniczkowaniu
$p_n'=-\lambda a(1-a)^{n-1}+\lambda(n-1)a^2(1-a)^{n-2}$.
Prawa strona równania Kołmogorowa wynosi
$-n\lambda a(1-a)^{n-1}+(n-1)\lambda a(1-a)^{n-2}$,
co po wyłączeniu $\lambda a(1-a)^{n-2}$ sprowadza się do
$-n(1-a)+(n-1)=na-1$,
tak samo jak pierwsza pochodna.
Jednoznaczność rozwiązań wynika indukcyjnie: każde równanie
jest liniowym równaniem niejednorodnym z wcześniej znaną
funkcją $p_{n-1}$. Ponieważ kandydat daje rozkład
prawdopodobieństwa ($\sum_{n\ge1}a(1-a)^{n-1}=1$),
odzyskujemy prawo procesu. Średnia rozkładu
geometrycznego na $\{1,2,\ldots\}$ wynosi $1/a=e^{\lambda t}$.
\end{proof}

\begin{uwaga}[Dwa odmienne mechanizmy wzrostu]
W procesie Galtona--Watsona wszyscy osobnicy rozmnażają się
w kolejnych \emph{pokoleniach}, a liczebność zmienia się
skokowo w chwilach całkowitych. W procesie Yule'a
urodzenia pojawiają się w \emph{czasie ciągłym}, przy czym
intensywność łączna rośnie proporcjonalnie do bieżącej liczby
osobników. Oba procesy mają strukturę rozgałęziającą,
ale inne trajektorie, rozkłady i narzędzia obliczeniowe.
\end{uwaga}


\section{Jak szybko zanika prawdopodobieństwo przetrwania?}
\begin{stwierdzenie}[Ogon czasu wymarcia w przypadku podkrytycznym]
Niech $Z$ będzie procesem Galtona--Watsona rozpoczynającym się
od $Z_0=1$ ze średnią potomstwa $m\in[0,1)$.
Niech $\tau=\inf\{n\ge0:Z_n=0\}$.
Wtedy dla każdego $n\ge0$
\[
\Prob(\tau>n)=\Prob(Z_n>0)\le m^n,
\qquad \E\tau\le\frac{1}{1-m}.
\]
\end{stwierdzenie}
\begin{proof}
Ponieważ zero jest stanem pochłaniającym,
$\{\tau>n\}=\{Z_n>0\}$. Dla nieujemnej zmiennej
całkowitoliczbowej nierówność Markowa daje
$\Prob(Z_n>0)=\Prob(Z_n\ge1)\le\E Z_n=m^n$.
Wzór ogonowy dla zmiennej o wartościach w $\N_0\cup\{\infty\}$,
wynikający z twierdzenia Tonellego, ma postać
$\E\tau=\sum_{n=0}^\infty\Prob(\tau>n)$.
Sumując oszacowania, otrzymujemy deklarowaną granicę.
\end{proof}

\begin{twierdzenie}[Asymptotyka Kołmogorowa w przypadku krytycznym]
\label{thm:gw-critical-survival}
Niech $Z_0=1$, $\E\xi=1$ i
$0<\sigma^2=\Var(\xi)<\infty$.
Wtedy
\[
\Prob(Z_n>0)\sim\frac{2}{\sigma^2 n}
\qquad(n\to\infty).
\]
W szczególności czas wymarcia $\tau$ jest skończony
prawie na pewno, lecz $\E\tau=\infty$.
\end{twierdzenie}
\begin{proof}
Warunki $m=1$ i $\sigma^2>0$ wymuszają $p_0>0$:
gdyby $\xi\ge1$ p.n., średnia $1$ wymuszałaby
$\xi=1$ p.n. Z kryterium wymierania wynika
$q_n=\Prob(Z_n=0)\uparrow1$.
Połóżmy $a_n=1-q_n>0$; dodatniość wynika
z istnienia dodatniego prawdopodobieństwa
co najmniej jednego dziecka w każdym z $n$
pokoleń (można śledzić pierwszą linię potomną).
Mamy $a_{n+1}=1-f(1-a_n)$.

Dla $0<s<1$ różniczkowanie szeregu potęgowego
daje $f''(s)=\sum_{k\ge2} k(k-1)p_k s^{k-2}$.
Skończoność drugiego momentu daje jednostajną
zbieżność szeregu pochodnych do rzędu drugiego
na $[0,1]$ (każdy wyraz jest ograniczony przez
$k(k-1)p_k$ dla drugiej pochodnej).
W szczególności $f'(1)=1$ oraz
$f''(1)=\E[\xi(\xi-1)]=\sigma^2$,
gdzie pochodne w $1$ są lewostronne.
Wzór Taylora daje zatem, dla $u\downarrow0$,
\[
1-f(1-u)=u-\tfrac12\sigma^2u^2+o(u^2).
\]
Oznaczmy $c=\sigma^2/2>0$. Ponieważ $a_n\to0$,
mamy $a_{n+1}/a_n\to1$ oraz
\[
\frac1{a_{n+1}}-\frac1{a_n}
=\frac{a_n-a_{n+1}}{a_na_{n+1}}\longrightarrow c.
\]
Sumując różnice i stosując twierdzenie o średnich
arytmetycznych ciągu zbieżnego, dostajemy
$1/(na_n)\to c$. Zatem $na_n\to1/c$.
Ze wzoru ogonowego $\E\tau=\sum_{n\ge0}a_n$
i rozbieżności szeregu harmonicznego wynika
ostatnia teza.
\end{proof}

\begin{twierdzenie}[Twierdzenie Yagloma]
\label{thm:gw-yaglom}
Przy założeniach twierdzenia~\ref{thm:gw-critical-survival}
warunkowy rozkład $Z_n/n$ przy $Z_n>0$ zbiega
według rozkładu do rozkładu wykładniczego
o średniej $\sigma^2/2$, czyli intensywności
$2/\sigma^2$:
\[
\Law\left(\frac{Z_n}{n}\,\middle|\,Z_n>0\right)
\longrightarrow\mathrm{Exp}(2/\sigma^2).
\]
\end{twierdzenie}
\begin{proof}
Niech $c=\sigma^2/2$ i
$G(u)=1-f(1-u)$ dla $0<u\le1$.
Dla $0<u<1$ mamy $0<G(u)\le u$,
ponieważ $f(1-u)\ge1-u$ w przypadku krytycznym,
a $f(1-u)<1$ dzięki $\Prob(\xi>0)>0$.
Z rozwinięcia użytego w poprzednim dowodzie
wynika
\[
D(u):=\frac1{G(u)}-\frac1u\longrightarrow c
\quad(u\downarrow0).
\]
Ustalmy $\theta>0$ i dla każdego $n\ge1$ połóżmy
$b_{n,0}=1-e^{-\theta/n}$ oraz
$b_{n,j+1}=G(b_{n,j})$ dla $0\le j<n$.
Mamy $0<b_{n,j}\le b_{n,0}\to0$ jednostajnie po $j$,
więc $\max_{0\le j<n}|D(b_{n,j})-c|\to0$.
Sumowanie różnic odwrotności daje
\[
\frac1{n b_{n,n}}
=\frac1{n b_{n,0}}+
\frac1n\sum_{j=0}^{n-1}D(b_{n,j})
\longrightarrow \frac1\theta+c .
\]
Zatem $n b_{n,n}\to\theta/(1+c\theta)$.
Ponieważ $b_{n,n}=1-f^{\circ n}(e^{-\theta/n})$
oraz $a_n=1-f^{\circ n}(0)$,
warunkowa transformata Laplace'a wynosi
\[
\E[e^{-\theta Z_n/n}\mid Z_n>0]
=\frac{f^{\circ n}(e^{-\theta/n})-f^{\circ n}(0)}
{a_n}
=1-\frac{b_{n,n}}{a_n}
\longrightarrow \frac1{1+c\theta}.
\]
Dla $\theta=0$ równość transformat jest automatyczna.
Ostatnia funkcja jest transformatą rozkładu
$\mathrm{Exp}(1/c)$.

Dla ścisłości wniosek o zbieżności rozkładów
można otrzymać bez dodatkowego twierdzenia
o transformatach. Średnie warunkowe
$\E[Z_n/n\mid Z_n>0]=1/(na_n)$ są ograniczone,
więc nierówność Markowa daje ciasność praw
na $[0,\infty)$. Każdy podciąg ma słabo zbieżny
podciąg (twierdzenie o wyborze Helly'ego dla
dystrybuant, z ciasnością). Jego granica ma
wskazaną transformatę, gdyż $x\mapsto e^{-\theta x}$
jest ograniczona i ciągła.
Jednoznaczność transformaty wynika już z wartości
dla $\theta=k\in\N_0$: po odwzorowaniu $x\mapsto e^{-x}$
otrzymujemy wszystkie momenty miary na $[0,1]$.
Wielomiany są jednostajnie gęste w $C([0,1])$
na mocy twierdzenia Weierstrassa, więc momenty
wyznaczają miarę, a odwzorowanie $x\mapsto e^{-x}$
jest różnowartościowe z borelowską odwrotnością
na $(0,1]$. Wszystkie granice podciągowe są
zatem identyczne, co kończy dowód.
\end{proof}

\begin{wniosek}[Średnia i wariancja pod warunkiem przetrwania]
Przy tych samych założeniach
\[
\E[Z_n/n\mid Z_n>0]\longrightarrow c,\qquad
\Var(Z_n/n\mid Z_n>0)\longrightarrow c^2,
\quad c=\sigma^2/2.
\]
\end{wniosek}
\begin{proof}
Ponieważ $\E Z_n=1$ oraz $\E Z_n^2=1+n\sigma^2$,
a oba momenty nie mają wkładu ze zdarzenia $Z_n=0$,
mamy
\[
\E[Z_n/n\mid Z_n>0]=\frac1{na_n},\qquad
\E[(Z_n/n)^2\mid Z_n>0]=\frac{1+n\sigma^2}{n^2a_n}.
\]
Asymptotyka $na_n\to1/c$ daje granice $c$ oraz
$\sigma^2c=2c^2$. Po odjęciu kwadratu średniej
otrzymujemy granicę wariancji.
\end{proof}

\begin{przyklad}[Krytyczne potomstwo zero lub dwa]
Dla $\Prob(\xi=0)=\Prob(\xi=2)=1/2$ mamy
$m=1$, $\sigma^2=1$ i
$a_{n+1}=a_n-a_n^2/2$, $a_0=1$.
Wymarcie następuje prawie na pewno, lecz
$\Prob(Z_n>0)\sim2/n$.
Warunkowo względem przetrwania liczebność jest
rzędu $n$, a $Z_n/n$ ma graniczny rozkład
wykładniczy o średniej $1/2$.
Nie ma sprzeczności z $\E Z_n=1$:
malejąca szansa przetrwania równoważy wzrost
średniej liczebności w ocalałych populacjach.
\end{przyklad}

\section{Urodzenia i śmierci w czasie ciągłym}
\begin{twierdzenie}[Prawdopodobieństwo wymarcia liniowego procesu urodzin i śmierci]
Niech $\lambda,\delta>0$. Rozważmy minimalny łańcuch Markowa
$Z$ na $\N_0$, rozpoczynający się od $Z_0=1$, z intensywnościami
$q_{n,n+1}=\lambda n$, $q_{n,n-1}=\delta n$ dla $n\ge1$,
$q_{nn}=-(\lambda+\delta)n$, zaś $0$ jest stanem pochłaniającym.
Wtedy proces jest nieeksplozywny, a prawdopodobieństwo
ostatecznego wymarcia wynosi
\[
q=\begin{cases}1,&\lambda\le\delta,\\
\delta/\lambda,&\lambda>\delta.
\end{cases}
\]
\end{twierdzenie}
\begin{proof}
Najpierw uzasadnimy brak eksplozji bez odwoływania się do
nieformalnego sprzężenia. Dla $K\ge2$ niech
$\tau_K=\inf\{t\ge0:Z_t\ge K\}$. Proces zatrzymany
$Z_{t\wedge\tau_K}$ przyjmuje wartości w skończonym
zbiorze $\{0,1,\ldots,K\}$. Dla funkcji $f(n)=n$
jego generator przed zatrzymaniem wynosi
$Qf(n)=(\lambda-\delta)n\le\lambda n$.
Wzór Dynkina (twierdzenie~\ref{thm:dynkin})
dla skończonego łańcucha z pochłaniającymi stanami
$0$ i $K$, który ma prawo procesu do zatrzymania, daje dla $u_K(t)=\E Z_{t\wedge\tau_K}$
\[
u_K(t)\le1+\lambda\int_0^tu_K(s)\,\mathrm d\lambda_1(s).
\]
Nierówność Gronwalla prowadzi do
$u_K(t)\le e^{\lambda t}$. Ponieważ
$Z_{t\wedge\tau_K}=K$ na $\{\tau_K\le t\}$,
$\Prob(\tau_K\le t)\le e^{\lambda t}/K$.
Eksplozja w skończonym czasie wymaga osiągnięcia
każdego poziomu $K$ poza zbiorem zerowym:
jeżeli łańcuch pozostaje w skończonym zbiorze stanów,
jego intensywności są ograniczone, a suma czasów
oczekiwania przy nieskończenie wielu skokach
jest nieskończona p.n. (twierdzenie~\ref{thm:bounded-rates},
lub sprzężenie z niezależnymi odstępami wykładniczymi
o wspólnej maksymalnej intensywności).
Ponieważ każdy skok zmienia liczebność o jeden,
ucieczka z każdego skończonego zbioru wymusza
osiągnięcie każdego $K$. Dla każdego $t<\infty$
$\Prob(\zeta\le t)\le\inf_{K\ge2}e^{\lambda t}/K=0$,
gdzie $\zeta$ oznacza czas eksplozji. Proces jest zatem
nieeksplozywny.

Rozważmy pierwsze zdarzenie po starcie od jednego osobnika.
Łączna intensywność wynosi $\lambda+\delta$, a jego rodzaj
jest niezależny od czasu oczekiwania: z prawdopodobieństwem
$\delta/(\lambda+\delta)$ osobnik umiera, z prawdopodobieństwem
$\lambda/(\lambda+\delta)$ rodzi drugiego osobnika.
Konstrukcję przez niezależne zegary uzasadniamy
następująco: każdemu osobnikowi przypisujemy niezależny
czas wykładniczy o intensywności $\lambda+\delta$
i niezależny typ zdarzenia (śmierć lub rozdzielenie
na dwóch osobników). Minimum $n$ zegarów ma
intensywność $n(\lambda+\delta)$, a typ zdarzenia
ma podane prawdopodobieństwa. Własność braku pamięci
daje po zdarzeniu niezależne nowe zegary pozostałych
osobników, więc konstrukcja ma dokładnie intensywności
i czasy przebywania z definicji minimalnego łańcucha.
Brak eksplozji pozwala prowadzić ją dla wszystkich czasów.
Po rozdzieleniu mamy dwa niezależne drzewa potomne,
każde o takim samym prawdopodobieństwie wymarcia $q$.
Wobec tego
\[
q=\frac{\delta}{\lambda+\delta}
+\frac{\lambda}{\lambda+\delta}q^2.
\]
Rozwiązaniami są $1$ i $\delta/\lambda$.
Uzasadnijmy wybór mniejszego pierwiastka. Każdy żyjący
osobnik ma niezależny czas do pierwszego zdarzenia
o rozkładzie $\mathrm{Exp}(\lambda+\delta)$; wtedy
z prawdopodobieństwem $\delta/(\lambda+\delta)$
umiera bez potomków, a z prawdopodobieństwem
$\lambda/(\lambda+\delta)$ rodzi nowego osobnika.
W drugim przypadku można potraktować dotychczasowego
osobnika (z uwagi na własność braku pamięci) i nowego
jako \emph{dwie niezależne linie potomne}. Drzewo
zdarzeń jest więc drzewem Galtona--Watsona z liczbą
dzieci $0$ lub $2$ i funkcją tworzącą
$f(s)=\delta/(\lambda+\delta)+\lambda s^2/(\lambda+\delta)$.
Zdarzenie, że drzewo ma skończoną liczbę węzłów,
jest sumą zdarzeń wymarcia do skończonej głębokości;
jego prawdopodobieństwo jest zatem granicą iteracji
$f^{\circ n}(0)$, czyli najmniejszym punktem stałym $f$.
Jeżeli drzewo jest skończone, wszystkie związane z nim
czasy oczekiwania są skończone prawie na pewno,
a więc proces rzeczywiście osiąga stan $0$ w skończonym
czasie. Odwrotnie, jeśli osiąga $0$ w skończonym czasie,
to brak eksplozji oznacza skończoną liczbę zdarzeń do tej
chwili, a zatem skończone drzewo genealogiczne.
Wymarcie procesu jest więc dokładnie wymarciem tego drzewa.
Stosując kryterium z wcześniejszej części rozdziału,
wybieramy najmniejszy pierwiastek: $q=1$ przy
$\lambda\le\delta$ i $q=\delta/\lambda$ przy
$\lambda>\delta$.
\end{proof}

\begin{zadanie}[Próg przeżycia populacji]
Niech każda żyjąca cząstka rodzi nową z intensywnością
$\lambda=3$ i ginie z intensywnością $\delta=2$,
niezależnie od pozostałych. Rozpoczynamy od jednej cząstki.
Jakie jest prawdopodobieństwo, że populacja kiedyś wymrze?

\emph{Rozwiązanie.} Pierwsze zdarzenie jest śmiercią
z prawdopodobieństwem $2/5$ lub narodzinami z prawdopodobieństwem
$3/5$. Stąd $q=2/5+(3/5)q^2$, więc
$(3q-2)(q-1)=0$.
Wybieramy mniejsze rozwiązanie: $q=2/3$.
Średni dodatni przyrost liczby osobników nie gwarantuje
przetrwania konkretnego drzewa genealogicznego.
\end{zadanie}


\begin{zadanie}[Wzrost procesu Yule'a]
Niech $Z$ będzie procesem Yule'a o intensywności $\lambda>0$
z jednym osobnikiem początkowym. Wyznaczyć
$\Prob(Z_t\ge k)$, $\Var(Z_t)$ i
$\Prob(Z_t\ge3\mid Z_t\ge2)$ dla $t>0$, $k\ge1$.

\emph{Rozwiązanie.} Wprowadzamy $a=e^{-\lambda t}$.
Z udowodnionego wzoru $\Prob(Z_t=n)=a(1-a)^{n-1}$
otrzymujemy przez zsumowanie szeregu geometrycznego
$\Prob(Z_t\ge k)=(1-a)^{k-1}$.
Dla rozkładu geometrycznego na $\{1,2,\dots\}$
$\E Z_t=1/a$ i $\Var(Z_t)=(1-a)/a^2$, a więc
$\Var(Z_t)=e^{2\lambda t}-e^{\lambda t}$.
Wreszcie
\[
\Prob(Z_t\ge3\mid Z_t\ge2)
=\frac{(1-a)^2}{1-a}=1-e^{-\lambda t}.
\]
To przykład odrębnej własności braku pamięci
\emph{rozkładu geometrycznego liczebności},
a nie wykładniczego rozkładu czasów między skokami.
\end{zadanie}

\section{Podsumowanie i kierunek dalszej teorii}
Łańcuch Markowa opisuje przejścia między stanami;
proces odnowy opisuje cykle; proces rozgałęziający
opisuje losową liczbę \emph{potomków}, którzy sami
stają się źródłem kolejnych losowań. Funkcje tworzące
zamieniają złożoną strukturę drzew na iterację funkcji,
a martyngały pozwalają śledzić zmienność populacji
w kolejnych pokoleniach. Dalszym krokiem będą martyngały w czasie ciągłym, lokalizacja,
zatrzymywanie i wariacja kwadratowa. Całka Itô i równania
stochastyczne zostaną skonstruowane dopiero w tomie III.


% ================================================================
% ROZDZIAŁ 11. MARTYNGAŁY W CZASIE CIĄGŁYM
% ================================================================
